\documentclass[twocolumn]{article}
\usepackage[utf8]{inputenc}
\usepackage{blindtext}
\usepackage[paperheight=252mm,paperwidth=174mm,margin=1mm,heightrounded]{geometry}
\usepackage{ulem}
\usepackage{array}
\usepackage{listings}
\usepackage{fvextra}
\usepackage{amsmath}
\usepackage{csquotes}

\usepackage{minted}
%\BeforeBeginEnvironment{minted}{}
%\AfterEndEnvironment{minted}{}

\usepackage{tikz}
\usetikzlibrary{shapes.geometric, arrows}
\usetikzlibrary{positioning}

%% Common TikZ libraries
\usetikzlibrary{calc}

%% Custom TikZ addons
\usetikzlibrary{crypto.symbols}
\tikzset{shadows=no}        % Option: add shadows to XOR, ADD, etc.

\usepackage{color}
\definecolor{light-gray}{rgb}{0.95,0.95,0.95}
\setminted{bgcolor=light-gray}  % this line causes the problem

\setlength{\parindent}{0mm}
\setlength{\parskip}{0mm}

\tikzstyle{string} = [rectangle, rounded corners, inner sep=2mm, text centered, draw=black, fill=red!30]
\tikzstyle{bytes} = [rectangle, rounded corners, inner sep=2mm, text width=, text centered, draw=black, fill=blue!30]
\tikzstyle{int} = [rectangle, rounded corners, inner sep=2mm, text width=,  text centered, draw=black, fill=green!30]
\tikzstyle{arrow} = [thick,->,>=stealth]

\begin{document}

\title{Type-Guided LLVM Obfuscation}
\date{}
%\maketitle

\section*{Type-Guided LLVM Obfuscation}
\vspace*{-0.5\baselineskip}

LLVM passes are popular for implementing obfuscations. Techniques such as control-flow flattening and mixed boolean arithmetic (MBA) are easily performed via transformation passes. A drawback with this is that since we are operating at the LLVM IR level, a lot of information, such as most  types, have already been erased at this stage. Types are powerful because they can encode invariants that can inform our obfuscation.

\vspace*{-0.5\baselineskip}
\subsection*{Invariant-based MBA}

As a simple example, consider a custom integer type that supports all the usual operations but internally stores the value $x$ as $2x$. If we have a function that takes an argument $x$ of this type, we could in principle replace every use of (the internal value of) this argument with $x + (x \& 1)$. For performance and simplicity reasons, most decompilers and deobfuscation tools will perform their analysis over all possible values and thus they are unable to simplify the expression back to $x$ since it can easily find a counter-example $1 + (1 \& 1) \neq 1$. This ``conditional MBA" uses non-local information to strengthen the obfuscation. While we can perform a decent amount of source-level obfuscation where we do have access to the type information, it would be even better if we could get the best of two worlds: rich information about invariants from the source-level with the fine-grained control at the LLVM IR level.

\vspace*{-0.5\baselineskip}
\subsection*{Clang annotations}

Luckily, at least for C and C++ code, we can use clang annotations to transfer information to the IR level. In C++, any constant expression can be used. For example, we can annotate a function parameter as follows:

\begin{minted}{c++}
int f([[clang::annotate("obf", 42)]] int a) {
\end{minted}

Which at the IR level will appear as wrapping the argument in a pointer which gets annotated with a reference to a structure containing the associated data.

\begin{minted}{llvm}
@.str = private unnamed_addr constant [4 x i8]
c"obf\00", section "llvm.metadata"
@.str.1 = private unnamed_addr constant
[6 x i8] c"a.cpp\00", section "llvm.metadata"
@.args = private unnamed_addr constant
{ i32 } { i32 42 }, section "llvm.metadata"
...
%2 = alloca i32, align 4
store i32 %0, ptr %2, align 4
call void @llvm.var.annotation.p0.p0(ptr %2,
ptr @.str, ptr @.str.1, i32 2, ptr @.args)
\end{minted}

Any sequence of constant values can be used. However, this annotation is expressed as a (useless) function call which is optimized away later in the LLVM pipeline. To make it more robust, we can transfer the data to a metadata node, which is independent of optimizations, and attach it to instructions reading from this pointer. 

To do this, we need to create a transformation pass and insert it early in the pipeline. In this pass, we iterate over instructions and find the call to the generated annotation:


\begin{minted}{c++}
// for each instruction:
auto *CI = dyn_cast<CallInst>(&I);

CI->getCalledFunction()
  ->getName()
  .starts_with("llvm.var.annotation")
\end{minted}

We can then extract the arguments to this function call and chase those pointers to the global data containing our annotation data, which can be seen in the IR snippet listed previously.

\begin{minted}{c++}
auto *GV = dyn_cast<GlobalVariable>(
    CI->getArgOperand(4)
    ->stripPointerCasts()
);
GV->getInitializer();
\end{minted}

Once we have the obfuscation data, we can also find all the places where the annotated pointer is used and create and attach a metadata node with this data:

\begin{minted}{c++}
CI->getOperand(0)
  ->stripPointerCasts()
  ->users()

// for each user
auto *CInt = 
ConstantInt::get(Type::getInt64Ty(Ctx), param);
auto *IntMD = 
ConstantAsMetadata::get(CInt);
MDNode *Node = MDNode::get(Ctx, IntMD);
Load->setMetadata("obfuscate", Node);

\end{minted}

Finally, at a later point, for example in a later pass, we can access this data when building obfuscations:

\begin{minted}{c++}
// for each instruction
if (MDNode *N = I.getMetadata("obfuscate")) {
   ...
\end{minted}

\vspace*{-0.5\baselineskip}
\subsection*{Conclusion}

This specific invariant and metadata here is simple but can be expanded. By combining this with C++ templates and/or macros, complex relationships can be expressed resulting in obfuscations with global relationships which makes local-only analysis ineffective. Maybe we want to use various divisibility properties, non-standard ways of representing integers or something even more exotic. What obfuscations can you come up with using this technique?

%\vfill\null
\end{document}
